\section{방법론}
\label{sec:method}

그림~\ref{fig:pipeline}은 본 연구의 절차를 네 블록으로 보인다.
데이터 수집(A)과 누수 없는 표본 구성(B)은 공개 텔레메트리를 팀--단계 예측 표본으로 바꾸고, 설계와 분석(C)은 세 연구 질문에 대응하는 모델을 구성하며, 평가(D)는 모든 설계를 같은 조건에서 비교한다.
이하의 각 절은 그림의 블록 순서를 따르며, 구현 세부는 익명 온라인 부록에 수록하였다.

\begin{figure*}[t]
\centering
\includegraphics[width=\textwidth]{fig_pipeline_clean.pdf}
\caption{데이터 파이프라인과 모델링·평가 절차}
\label{fig:pipeline}
\end{figure*}

%======================================================================
\subsection{데이터 수집}
\label{sec:method:data}
PUBG의 공개 텔레메트리는 경기에 참가한 모든 플레이어의 위치를 약 10초 간격으로 기록하고, 피해, 기절, 처치, 아이템, 안전구역 상태, 단계 전환과 같은 사건을 발생 시각과 함께 기록한다~\cite{PUBGAPI}.
본 연구는 공개 API의 표본 경기 목록을 하루 단위로 조회하여 스팀 서버의 경기를 수집하였으며, 수집한 경기는 모두 2026년 6월의 안전구역 규칙 개편~\cite{PUBGBlueZone2026} 이전의 같은 규칙을 따른다.
경기 중 모든 플레이어의 위치를 이 해상도로 담은 공개 자료는 드물며, 기존 배틀로얄 예측 연구가 사용한 경기 내 기록도 공개되지 않았다~\cite{Sagong25a,KRAFTON}.
분석에 사용한 경기와 표본의 규모는 표~\ref{tab:cohorts}에 정리하였다.

%======================================================================
\subsection{누수 없는 표본 구성}
\label{sec:method:samples}

\subsubsection{시간축, 예측 시점, 입력 절단}
\rev{모든 기록의 타임스탬프를 경기 시작 사건으로부터의 경과 시간(초)으로 변환하여 하나의 시간축을 만들고, 그 위의 단계 전환 사건으로 단계를 정한다.
경기 $m$의 단계 $k$는 $k$번째 안전구역이 공개된 시점부터 그 구역으로의 축소가 끝날 때까지이며, 축소가 끝나 다음 구역이 공개되는 시각 $b_{m,k}$를 예측 시점으로 둔다.
모든 입력은 특징 절단 시각 $a_{m,k}=b_{m,k}-\Delta$ 이전의 기록만으로 만들며($\Delta=10$초), 이 제한은 팀 특징뿐 아니라 생존 판정, 팀 간 관계, attention 마스크에도 적용한다.
생존 여부는 $a_{m,k}$ 이전의 처치와 위치 기록으로 판정하며, 처치된 플레이어는 이후 체력이 남은 위치 기록이 나타날 때에만 재투입된 것으로 본다.
무작위로 고른 경기에서 $a_{m,k}$ 이후의 기록을 모두 지우고 입력을 다시 만들어도 원래 입력과 일치함을 확인하였다.}

\subsubsection{사건 기반 라벨과 예측 대상}
\rev{기절은 최종 처치와 구분하며, 경기 $m$에서 팀 $i$의 탈락 시각 $T_{m,i}$는 재투입되지 않은 마지막 팀원이 처치된 시각이다. 우승 팀은 탈락하지 않으므로 $T_{m,i}=\infty$로 두고, 경기 종료 시각 $C_m$에서 관측을 종료한다.
단계 $k$의 위험집합은 예측 시점에 생존한 팀의 집합 $\mathcal{R}_{m,k}=\{i: T_{m,i}>b_{m,k}\}$이고, 라벨은 $Y_{m,i,k}=\mathbb{1}[T_{m,i}\le b_{m,k+1}]$이다. 여기서 $\mathbb{1}[\cdot]$은 조건이 참이면 1, 거짓이면 0인 지시함수이며, 단계 $k+1$이 경기의 마지막 단계이면 $b_{m,k+1}=C_m$으로 둔다.
분석 단위는 $k=1,\dots,7$의 단계 경계에서 만든 팀--단계 표본 $(m,i,k)$이며, 솔로 경기에서는 곧 플레이어의 표본이다.
예측 대상은 다음 단계 탈락 위험}
\begin{equation}
\lambda_{m,i,k}=\Pr\!\left(Y_{m,i,k}=1 \mid i\in\mathcal{R}_{m,k},\,\mathcal{H}_{m}(a_{m,k})\right)
\label{eq:hazard}
\end{equation}
\rev{으로, 단계 경계를 landmark로 하는 이산시간 hazard이다~\cite{SingerWillett93,vanHouwelingen07}. 여기서 $\mathcal{H}_{m}(t)$는 경기 $m$에서 시각 $t$ 이전에 기록된 텔레메트리 전체이다.}

\subsubsection{코호트와 분할}
\label{sec:method:cohort}
분석 대상은 일반 경기와 랭크 경기로 한정하고, 튜토리얼과 훈련장 경기, 계정 대부분이 AI인 캐주얼 모드, 비플레이어 전투 계정이 다수 기록되는 Vikendi 맵의 경기는 제외하였다.
경기는 맵, 시점, 팀 규모의 조합으로 코호트를 나누어 충분한 경기가 있는 12개 조합을 사용하였고(표~\ref{tab:cohorts}), 경기 단위로 학습, 검증, 시험 집합에 8:1:1로 고정 분할하였다.
일반 경기에도 AI 계정이 섞여 있으므로 각 팀의 AI 계정 비율을 입력에 포함하였다.

\begin{table}[t]
\centering
\caption{코호트별 경기 수, 팀--단계 표본 수, 다음 단계 탈락률}
\label{tab:cohorts}
\footnotesize
\setlength{\tabcolsep}{3pt}
\begin{tabular}{lccrrrr}
\hline
맵 & 시점 & 인원 & 경기 & 학습/검증/시험 & 표본 & 탈락률(\%) \\
\hline
Erangel & 3인칭 & 1 & 224 & 189/15/20 & 31,740 & 35.1 \\
Erangel & 3인칭 & 2 & 656 & 526/65/65 & 59,170 & 29.5 \\
Taego & 3인칭 & 2 & 573 & 459/57/57 & 47,223 & 29.2 \\
Erangel & 1인칭 & 4 & 430 & 344/43/43 & 21,849 & 27.8 \\
Erangel & 3인칭 & 4 & 3,283 & 2,638/321/324 & 196,442 & 25.7 \\
Karakin & 3인칭 & 4 & 333 & 267/33/33 & 12,669 & 32.1 \\
Miramar & 3인칭 & 4 & 1,543 & 1,235/154/154 & 92,279 & 24.6 \\
Paramo & 3인칭 & 4 & 333 & 267/33/33 & 9,963 & 30.5 \\
Rondo & 3인칭 & 4 & 1,039 & 833/103/103 & 53,683 & 25.2 \\
Sanhok & 3인칭 & 4 & 940 & 752/94/94 & 48,793 & 27.6 \\
Taego & 1인칭 & 4 & 435 & 349/43/43 & 23,390 & 30.0 \\
Taego & 3인칭 & 4 & 4,159 & 3,329/415/415 & 220,565 & 24.9 \\
\hline
합계 & & & 13,948 & 11,188/1,376/1,384 & 817,766 & 26.6 \\
\hline
\end{tabular}
\\[2pt]
\parbox{\columnwidth}{\scriptsize 2026년 3월 26일--6월 2일에 생성된 경기로, 텔레메트리 기록은 모두 약 3억 2,240만 건(위치 약 9,643만, 사건 약 2억 2,597만)이다. 표본은 위험집합에 속한 팀--단계 표본이며, 탈락률은 그중 다음 단계 안에 탈락한 비율이다.}
\end{table}

%======================================================================
\subsection{설계와 분석}
\label{sec:method:design}
그림~\ref{fig:pipeline}의 C와 같이, 서식지 맥락 표현은 RQ1의 해석 모델과 RQ2의 본 연구 설계에 입력되고, 두 기존 설계의 입력은 RQ2와 RQ3의 비교에 사용된다.

\subsubsection{서식지 맥락 표현}
\label{sec:method:representation}
각 팀의 상태를 서식지 축소 관점의 네 구성개념으로 정의한다.
구역 경계에 대한 노출은 팀 중심과 안전구역 경계 사이의 부호 있는 거리를 구역 밖 깊이와 구역 안쪽 깊이로 나누어 측정하고, 이동 부담은 공개된 안전구역에 도달하는 데 필요한 속도와 실제 이동 속도의 차이로 측정한다~\cite{Nathan08}.
주변 경쟁 밀도는 팀 주변 200\,m 안의 다른 팀 수, 남은 팀원 수는 절단 시각에 생존이 관측된 팀원 수이다.
이 구성개념은 경기 맥락, 팀 상태, 이동, 자원, 공간에 관한 특징과 함께 팀--단계 입력을 이루며, 같은 단계의 팀 사이에는 근접, 교전, 공동 압박의 세 관계를 정의한다.

\subsubsection{RQ1: 생태 hazard 모델}
\label{sec:method:rq1}
RQ1에 답하기 위해 구성개념만으로 위험을 설명하는 로지스틱 회귀
\begin{equation}
\mathrm{logit}\,\lambda_{m,i,k}=\alpha_k+\textstyle\sum_{c\in\mathcal{C}}\beta_c z_{c,m,i,k}
\label{eq:ehm}
\end{equation}
\rev{를 코호트마다 추정한다. 여기서 $\mathrm{logit}\,p=\log\{p/(1-p)\}$이고, $\alpha_k$는 단계 $k$의 기저위험을 나타내는 절편이다. $\mathcal{C}$는 구역 밖 깊이, 구역 안쪽 깊이, 이동 부족분, 주변 경쟁 밀도, 남은 팀원 수의 다섯 변수와 통제 변수인 절단 시각의 생존 팀 수의 로그로 이루어지며, 솔로 코호트에서는 남은 팀원 수를 뺀다. $z_{c,m,i,k}$는 변수 $c$를 그 코호트의 학습 집합 평균과 표준편차로 표준화한 값, $\beta_c$는 그 계수이다.}
표준오차는 경기 단위로 군집화하고, 코호트별 계수는 확률효과 모형으로 통합하여 방향, 크기, 이질성, 부호가 같은 코호트의 수를 보고한다~\cite{DerSimonian86}.
단계에 따른 변화는 구성개념과 단계 구간(초반 1--2, 중반 3--5, 후반 6--7)의 상호작용으로 검정한다.

\subsubsection{RQ2: 본 연구 설계와 기존 설계}
\label{sec:method:rq2}
본 연구의 설계는 위험을 단계별 기저위험과 팀 상태의 함수
\begin{equation}
\mathrm{logit}\,\lambda_{m,i,k}=\alpha_k+f(\mathbf{x}_{m,i,k},\mathcal{G}_{m,k})
\label{eq:logit}
\end{equation}
\rev{로 둔다. 여기서 $\mathbf{x}_{m,i,k}$는 팀 $i$의 서식지 맥락을 포함한 입력 특징 벡터, $\mathcal{G}_{m,k}$는 $a_{m,k}$에 생존한 팀들 사이의 세 관계를 담은 그래프, $f$는 학습되는 함수이다. 모델은 학습 집합의 모든 팀--단계 표본에 대한 교차엔트로피 $-\sum[Y_{m,i,k}\log\lambda_{m,i,k}+(1-Y_{m,i,k})\log(1-\lambda_{m,i,k})]$를 최소화한다.}
학습기는 단계 번호를 입력으로 받는 그래디언트 부스팅 결정나무(GBDT)~\cite{Friedman01,Ke17}와, 구조 요소의 기여를 분해하기 위한 참조 모델인 PTR-Former를 사용한다.
PTR-Former는 같은 팀의 단계 이력에 대한 인과적 attention(단계 축)과 같은 단계의 생존 팀에 대한 attention(팀 축)을 번갈아 적용하고~\cite{Vaswani17}, 팀 축 attention에 관계 유형별 편향을 더한다~\cite{Ying21}.

\rev{KRAFTON의 예측기는 Cox 부분우도의 항 $\exp(s_i)/\sum_{j\in\mathcal{R}}\exp(s_j)$에 가중치를 둔 목적을 학습한다~\cite{Cox72}. 여기서 $s_i$는 한 시점에서 팀 $i$의 점수, $\mathcal{R}$은 그 시점에 생존한 팀의 집합이다.
모든 팀의 점수에 같은 상수를 더해도 이 항은 변하지 않으므로, 이 목적은 식~\eqref{eq:logit}의 $\alpha_k$와 같은 단계별 위험 수준을 식별하지 못한다.
따라서 상대위험 점수에 단계별 기저위험 $\gamma_k$를 결합한 확률 $1-\exp\{-\exp(\gamma_k+s_{m,i,k})\}$도 함께 비교한다. 여기서 $s_{m,i,k}$는 예측 시점 $b_{m,k}$에서 팀 $i$의 점수이고, $\gamma_k$는 위험집합의 관측 탈락과 점수로부터 추정한 단계 $k$의 기저위험이다~\cite{PrenticeGloeckler78}.}
두 기존 설계는 원래의 입력, 학습 목적, 학습 표집을 유지한 채 같은 분할과 같은 위험집합에서 평가하였으며, 본 연구 설계와의 차이와 적용 방식은 표~\ref{tab:compare}에 정리하였다.
적용판이 불리하게 구성되지 않았는지 확인하기 위해, 각 적용판을 원래 과제의 지표(우승 팀 정확도와 탈락 순서 C-index, 솔로 코호트의 개인 생존 분류)로도 평가하고 본 연구의 설계와 비교한다.

\begin{table*}[t]
\centering
\caption{기존 설계와 본 연구 설계의 비교}
\label{tab:compare}
\footnotesize
\renewcommand{\arraystretch}{1.2}
\begin{tabular}{p{2.3cm}p{4.3cm}p{4.6cm}p{4.6cm}}
\hline
 & Sagong 등~\cite{Sagong25a} & KRAFTON~\cite{KRAFTON} & 본 연구 \\
\hline
분석 단위 & 플레이어 & 팀 & 팀(솔로 경기에서는 개인) \\
예측 대상 & 다음 반 단계의 생존 여부 & 탈락 순서의 상대위험(우승확률로 변환) & \textbf{다음 단계 안에 탈락할 조건부 확률} \\
예측 시점 & 반 단계 경계 & 경기 중 여러 시점 & 단계 경계(다음 안전구역 공개 시점) \\
단계별 위험 수준 & 플레이 유형별 전이 생존율 & 표현하지 않음 & \textbf{단계별 기저위험으로 명시} \\
입력 & 단계별 활동·교전 점수, 플레이 유형 & 팀의 행동·상태 지표, 팀원 위치 & \textbf{서식지 맥락}, 팀 간 관계 \\
팀 간 상호작용 & 다루지 않음 & 팀 간 attention & 팀 축·단계 축 attention과 관계 편향 \\
학습 목적 & 생존 분류 & 가중 Cox 부분우도 & \textbf{다음 단계 탈락의 교차엔트로피} \\
\hline
본 연구의 적용 & 2·4인 팀은 팀 단위로 집계, 예측 시점을 단계 경계로 통일 & 입력을 본 연구의 텔레메트리로 재구성, 출력을 확률로 보정~\cite{Platt99} & --- \\
\hline
\end{tabular}
\end{table*}

\subsubsection{RQ3: 단일 요인 분해}
\label{sec:method:decomposition}
RQ3에서는 다른 조건을 고정한 채 한 번에 하나의 요인만 바꾸는 단계를 이어 비교한다.
입력 표현은 학습기를 GBDT로 고정하고 Sagong 등의 입력에 서식지 맥락을 더한 뒤(a1) 나머지 입력을 본 연구의 입력으로 바꾸어(a2), 그리고 KRAFTON 입력을 본 연구의 입력으로 바꾸어(b) 측정한다.
학습 목적과 기저위험은 KRAFTON의 구조에 본 연구의 입력을 넣고 Cox 목적으로 학습한 모델에서 출발하여, 단계별 기저위험의 결합(c1), 다음 단계 교차엔트로피로의 목적 교체(c2), 명시적 단계 절편의 추가(d)로 측정한다.
구조는 KRAFTON 구조를 팀 축만 남긴 PTR-Former로 바꾸는 아키텍처 단계(e)와 단계 축과 관계 편향을 더하는 단계(f)로 측정하며, 이 두 단계에서는 순위 보조 손실을 쓰지 않는다.
마지막으로 GBDT에 KRAFTON 입력을 더하여 두 표현의 보완성을 확인한다(g).
이와 별도로 단계 축과 팀 축의 2$\times$2 요인 실험, 관계를 섞은 대조 모델, 구역 객체 토큰으로 구성한 다섯 비교는 실험 전에 지정하였다.

%======================================================================
\subsection{평가}
\label{sec:method:evaluation}
그림~\ref{fig:pipeline}의 D와 같이, 모든 설계의 예측은 같은 시험 표본에서 지표로 요약하고 코호트별 차이를 통합한다.

\subsubsection{평가 지표}
주 지표는 시험 위험집합 전체의 AUC이며, 단계마다 공통으로 더해지는 위험 수준을 제외한 판별력을 보기 위해 단계 안 AUC와 경기--단계 내부 AUC를 함께 보고한다.
확률 예측의 정확도는 로그 손실과 Brier 점수로 평가한다. \rev{단계별 위험 수준의 표현은 단계별 보정 오차 $\sum_{k}(n_k/n)\,|\bar{p}_k-\bar{y}_k|$로 평가하며, $n_k$는 단계 $k$의 시험 표본 수, $n=\sum_k n_k$, $\bar{p}_k$와 $\bar{y}_k$는 그 표본의 평균 예측확률과 관측 탈락률이다.}

\subsubsection{통계적 비교}
두 설계의 차이는 코호트마다 경기 단위의 쌍대 bootstrap(2,000회)으로 신뢰구간을 구하고~\cite{Efron93}, DerSimonian--Laird 확률효과 모형으로 통합하며~\cite{DerSimonian86}, REML 추정과 Hartung--Knapp 구간으로 민감도를 확인한다~\cite{Viechtbauer05,HartungKnapp01}.
AUC 차이는 의미 있는 최소 효과 크기 0.005를 기준으로 개선, 저하, 실질적 동등, 불확정으로 분류하고, 사전 지정한 다섯 비교에는 Holm 보정을 적용하며~\cite{Holm79}, 나머지는 탐색적 분석으로 보고한다.
강건성은 가장 큰 두 코호트에서 학습 초기화, 두 배의 학습 예산, 학습 경기 수를 바꾼 학습 곡선으로 확인한다.
